\subsection{string\_hash}\label{string_hash}
Das \verb|string_hash| Modul ist eine Hash-Map, die kurze Zeichenketten auf
einen \verb|int8_t| Wert abbildet. Ein Schlüssel ist so lang wie das Namensfeld
\verb|line| einer Einzelmessung der Zeitmessung (\verb|time_meas_t|,
Abschnitt \ref{time}). Typische Anwendung ist das Nachschlagen eines Handles anhand
seines Namens.

Die Bibliothek kommt ohne \verb|malloc| aus: Alle Knoten liegen in einem Pool
fester Grösse innerhalb der Map-Struktur. Wo die Map liegt (global, als
\verb|static| oder auf dem Stack), bestimmt der Anwender. Speicherbedarf und
Laufzeit sind damit deterministisch, und mehrere unabhängige Maps sind
möglich.

\subsubsection*{Konfiguration}
Die Grössen lassen sich vor dem Einbinden von \verb|string_hash.h| (oder
per Compiler-Option \verb|-D|) überschreiben (Listing \ref{sh_defines}).

\begin{listing}[H]
\begin{minted}{C}
#define SH_SIZE        100  // max. Anzahl Eintraege (Knoten-Pool)
#define SH_BUCKET_CNT   37  // Anzahl Buckets, Primzahl != 31
#define SH_KEY_LEN (TIME_MEAS_CHAR_PER_LINE + 1)  // inkl. '\0'
\end{minted}
\caption{Konfiguration von \texttt{string\_hash}}
\label{sh_defines}
\end{listing}

Ein Schlüssel darf höchstens \verb|SH_KEY_LEN - 1| Zeichen lang sein (bei
\verb|TIME_MEAS_CHAR_PER_LINE = 20| also 20 Zeichen); längere werden
abgewiesen und nicht abgeschnitten. \verb|SH_BUCKET_CNT| darf kein
Vielfaches des Hash-Multiplikators 31 sein, sonst fallen alle Zeichen
ausser dem letzten aus dem Hash und die Einträge verteilen sich schlecht.

Ein Knoten belegt auf dem Cortex-M4 28 Bytes, ein Bucket 4 Bytes. Mit den
Standardwerten braucht eine Map damit
$100 \cdot 28 + 37 \cdot 4 + 4 = 2952$\,Bytes $\approx 2{,}9$\,kB RAM.

\subsubsection*{Zustand}
Die Map wird durch \verb|sh_hash_map_t| beschrieben (Listing
\ref{sh_hash_map_t}). \verb|bucket| enthält pro Bucket den Kopf einer
verketteten Liste, \verb|node| ist der Knoten-Pool. Ein Knoten mit
\verb|used == false| ist frei. Kollisionen werden durch Verkettung
(\glqq separate chaining\grqq{}) aufgelöst; neue Knoten werden am Kopf der
Liste eingefügt.

\begin{listing}[H]
\begin{minted}{C}
typedef struct sh_node_s {
    char key[SH_KEY_LEN];    // Schluessel, '\0'-terminiert
    int8_t value;            // zugeordneter Wert
    bool used;               // Knoten belegt
    struct sh_node_s *next;  // naechster Knoten im Bucket
} sh_node_t;

typedef struct sh_hash_map_s {
    int16_t numOfElements;            // Anzahl Eintraege
    sh_node_t *bucket[SH_BUCKET_CNT]; // Listenkopf je Bucket
    sh_node_t node[SH_SIZE];          // Knoten-Pool
} sh_hash_map_t;
\end{minted}
\caption{Definition der \texttt{sh\_hash\_map\_t} Struktur}
\label{sh_hash_map_t}
\end{listing}

\subsubsection*{Hash-Funktion}
Der Bucket-Index ist ein Polynom-Hash nach dem Horner-Schema über die
Zeichen des Schlüssels, modulo der Anzahl Buckets:
\[
h = \left(\sum_{i=0}^{n-1} k_i \cdot 31^{\,n-1-i}\right) \bmod
\mathtt{SH\_BUCKET\_CNT}
\]
Gerechnet wird in \verb|uint32_t|. Ein Überlauf bei langen Schlüsseln ist
gewollt und harmlos, da das Ergebnis nur gleichmässig verteilt, aber nicht
exakt sein muss.

\subsubsection*{API}
Alle Funktionen prüfen ihre Pointer-Argumente und die Schlüssellänge.
Ungültige Argumente führen zu \verb|EM_ERR|.

{\sloppy\emergencystretch=3em
\begin{description}
\item[\texttt{sh\_init(mp)}] Setzt die Map zurück: alle Buckets leer, alle
  Knoten frei, \verb|numOfElements = 0|. Muss vor der ersten Verwendung
  aufgerufen werden. Rückgabe: \verb|em_msg| (\verb|EM_ERR| bei
  \verb|NULL|-Map, sonst \verb|EM_OK|).
\item[\texttt{sh\_insert(mp, key, value)}] Fügt \verb|key| mit \verb|value|
  ein. Der Schlüssel wird in den Knoten kopiert, der Aufrufer muss ihn also
  nicht aufbewahren. Existiert \verb|key| bereits, wird nur dessen Wert
  überschrieben. Rückgabe: \verb|em_msg| (\verb|EM_OK| nach Einfügen oder
  Überschreiben, \verb|EM_ERR| bei ungültigem Argument, zu langem Schlüssel
  oder vollem Pool).
\item[\texttt{sh\_search(mp, key, value)}] Sucht \verb|key| und schreibt den
  zugehörigen Wert nach \verb|*value|. Da jeder \verb|int8_t| ein gültiger
  Wert ist, wird \glqq nicht gefunden\grqq{} über den Rückgabewert gemeldet
  und nicht über einen reservierten Wert. Rückgabe: \verb|em_msg|
  (\verb|EM_OK|, wenn gefunden; \verb|EM_ERR|, wenn \verb|key| fehlt oder ein
  Argument ungültig ist, \verb|*value| bleibt dann unverändert).
\item[\texttt{sh\_delete(mp, key)}] Entfernt \verb|key| aus seiner Liste und
  gibt den Knoten an den Pool zurück. Rückgabe: \verb|em_msg| (\verb|EM_OK|
  nach dem Entfernen, \verb|EM_ERR|, wenn \verb|key| fehlt oder ein Argument
  ungültig ist).
\item[\texttt{sh\_hash\_function(key)}] Liefert den Bucket-Index von
  \verb|key|. Für die normale Verwendung nicht nötig, aber für Tests nützlich.
  Rückgabe: \verb|uint16_t| (Index im Bereich
  $0 \ldots \mathtt{SH\_BUCKET\_CNT} - 1$).
\end{description}
\par}

\subsubsection*{Beispiel}
\begin{listing}[H]
\begin{minted}{C}
static sh_hash_map_t names;
int8_t hdl;

sh_init(&names);
sh_insert(&names, "loop", 3);
if (sh_search(&names, "loop", &hdl) == EM_OK) {
    time_start(hdl, 0, NULL, NULL);
}
sh_delete(&names, "loop");
\end{minted}
\caption{Verwendung von \texttt{string\_hash}}
\label{sh_example}
\end{listing}

\subsubsection*{Grenzen}
Die Map ist nicht thread- bzw.\ interruptsicher; gleichzeitige Zugriffe
aus ISR und Hauptschleife müssen vom Aufrufer geschützt werden. Die Suche
nach einem freien Knoten ist linear im Pool ($O(\mathtt{SH\_SIZE})$), Suchen
und Löschen sind im Mittel $O(1)$.
